\section{Gambaran Umum Sistem Kerja Magang}
Selama pelaksanaan magang di PT Inti Utama Solusindo, terdapat beberapa sistem
dan alat kerja yang digunakan sehari-hari. Bagian ini menjelaskan alur kerja
secara umum, mulai dari orientasi awal hingga mekanisme pengelolaan tugas dan
evaluasi kinerja.

\subsection{Orientasi Karyawan Baru}
Hari pertama magang diisi dengan kegiatan Orientasi Karyawan Baru (OKB) yang
dilaksanakan di Ruang Meeting, Gedung \textit{Head Office} (HO) bersama
karyawan baru lainnya. Orientasi ini mencakup pengenalan budaya kerja, aturan,
dan kebijakan yang berlaku di perusahaan. Setelah itu, dilanjutkan dengan
aktivasi akun GitLab dan OpenProject, serta proses \textit{cloning} repositori
perusahaan untuk mempersiapkan \textit{tools} yang diperlukan untuk memulai
pekerjaan.

\subsection{Pengenalan Produk dan Environment}
Dalam proses pengenalan, dijelaskan bahwa saat ini PT Inti Utama Solusindo
sedang gencar melakukan \textit{development} untuk dua aplikasi yang ada
diberada dalam ekosistem Hospinet, di antaranya HealthyOne dan Canvasser.

\begin{figure}[H]
      \centering
      \includegraphics[width=10cm]{screen-healthyone-main.png}
      \caption{Tampilan Beranda HealthyOne}
\end{figure}

Pada tahap awal, diberikan akses versi \textit{development} dari aplikasi
tersebut untuk memahami alur fitur, struktur halaman, dan perilaku antarmuka
sebelum masuk ke tahap lebih lanjut. \textit{Environment staging} digunakan
sebagai tempat pengujian sebelum perubahan/fitur baru dipromosikan ke
\textit{environment production}.

\begin{figure}[H]
      \centering
      \includegraphics[width=10cm]{screen-canvasser-hunter-main.png}
      \caption{Tampilan Beranda Canvasser}
\end{figure}

\subsection{OpenProject}
Dalam pengelolaan tugas dan proyek, PT Inti Utama Solusindo menggunakan
OpenProject sebagai platform utama. OpenProject merupakan aplikasi open-source
project management yang mendukung pencatatan, distribusi, pemantauan, dan
evaluasi tugas secara terpusat. Melalui platform ini, mentor dapat memberikan
tugas secara langsung kepada assignee yang bertanggung jawab mengerjakan
pekerjaan tersebut. Apabila tidak terdapat tugas yang secara langsung
diberikan, mahasiswa juga dapat mengambil tugas yang masih terbuka dan belum
memiliki assignee, sehingga produktivitas kerja tetap terjaga.

\begin{figure}[H]
      \centering
      
\includegraphics[width=4cm]{logo-openproject.png}
      \caption{Logo OpenProject}
\end{figure}

\subsection{Siklus Task OpenProject}
\begin{enumerate}
      \item Perencanaan dan persetujuan
            \begin{itemize}[noitemsep, topsep=0pt]
                  \item \textbf{MJ Planning:} digunakan untuk menandai bahwa task masih berada pada tahap
                        perencanaan awal oleh tim terkait
                  \item \textbf{UI/UX Design:} dipakai saat rancangan tampilan dan alur penggunaan mulai disusun
                        sebagai dasar pengembangan task.
                  \item \textbf{PB Approved:} menandai bahwa task telah memperoleh persetujuan untuk masuk ke
                        tahap pengerjaan.
            \end{itemize}
      \item Pengerjaan oleh developer \begin{itemize}[noitemsep, topsep=0pt]
                  \item \textbf{IT To Do:} digunakan saat task sudah masuk ke daftar kerja developer, tetapi
                        belum mulai dikerjakan.
                  \item \textbf{IT On Progress:} dipakai ketika task sedang dikerjakan oleh developer
                  \item \textbf{IT Completed - Not Deployed:} menandai bahwa proses pengerjaan telah selesai,
                        tetapi hasilnya belum di-deploy.
            \end{itemize}
      \item Pengujian QA pada environment development/staging
            \begin{itemize}[noitemsep, topsep=0pt]
                  \item \textbf{QA To Do:} dipakai ketika task sudah masuk ke antrian pengujian oleh tim QA.
                  \item \textbf{QA In Testing Dev:} digunakan saat task sedang diuji pada environment
                        development.
                  \item \textbf{QA Ready for Staging:} menandai bahwa task telah siap dipindahkan ke environment
                        staging untuk pengujian lanjutan.
                  \item \textbf{Staging Deployed:} dipakai setelah hasil task berhasil di-deploy ke environment
                        staging.
                  \item \textbf{QA In Testing Staging:} digunakan saat task menjalani pengujian lanjutan pada
                        environment staging.
                  \item \textbf{QA Ready for Promote:} menandai bahwa task telah lolos pengujian di staging dan
                        siap dipromosikan ke production.
            \end{itemize}
      \item Promote ke production dan validasi akhir
            \begin{itemize}[noitemsep, topsep=0pt]
                  \item \textbf{QA Promoted:} dipakai setelah task berhasil dipindahkan dari staging ke
                        production.
                  \item \textbf{QA Failed Production:} dipakai ketika pada production masih ditemukan kendala
                        atau hasilnya belum sesuai harapan.
                  \item \textbf{QA Pass Production:} dipakai ketika task berjalan dengan baik pada production
                        sesuai hasil yang diharapkan.
                  \item \textbf{QA Non Issue:} digunakan ketika validasi di production tidak menemukan kendala
                        pada task tersebut.
                  \item \textbf{QA Done:} menandai bahwa seluruh proses pengujian dan validasi task telah
                        selesai.
            \end{itemize}
      \item Kemungkinan pembatalan
            \begin{itemize}[noitemsep, topsep=0pt]
                  \item \textbf{MJ Cancelled:} dipakai ketika task dibatalkan dan tidak dilanjutkan ke proses
                        berikutnya.
            \end{itemize}

\end{enumerate}
\subsection{Sistem KPI}
Untuk mendukung pengukuran produktivitas kerja, perusahaan menerapkan sistem
Key Performance Indicator (KPI) berbasis poin pada setiap task. Bobot poin yang
diberikan berkisar antara 5 hingga 30 poin, tergantung pada tingkat
kompleksitas dan beban kerja yang dibutuhkan. Dalam sistem ini, mahasiswa
magang ditargetkan untuk mencapai minimal 35 poin, sedangkan karyawan memiliki
target minimal 50 poin. Pencapaian KPI menjadi salah satu indikator evaluasi
performa selama periode kerja. Apabila target minimum tidak tercapai, hal
tersebut dapat memengaruhi penilaian kinerja magang. Sebaliknya, pencapaian
yang melampaui target dapat dikonversikan ke dalam bentuk insentif tambahan
sesuai kebijakan perusahaan.

\subsection{GitLab}
Selain OpenProject, perusahaan menggunakan GitLab untuk manajemen repository
dan kolaborasi pengembangan source code. GitLab berfungsi sebagai tempat
penyimpanan code, pengelolaan branch, dokumentasi perubahan, serta pengajuan
merge request.
\begin{figure}[H]
      \centering
      
\includegraphics[width=4cm]{logo-gitlab.png}
      \caption{Logo GitLab}
\end{figure}
\begin{figure}[H]
      \centering
      \includegraphics[width=10cm]{screen-gitlab-main.png}
      \caption{Tampilan GitLab}
\end{figure}

Setelah suatu task selesai dikerjakan, mahasiswa melakukan push perubahan ke
branch terkait, kemudian mengajukan merge request agar code dapat ditinjau
sebelum digabungkan ke branch utama, yaitu master. Proses ini membantu menjaga
konsistensi, keamanan, dan kualitas code dalam pengembangan aplikasi.
\begin{figure}[H]
      \centering
      \includegraphics[width=10cm]{screen-gitlab-mr.png}
      \caption{Tampilan Menu \textit{Merge Request} di GitLab}
\end{figure}

\begin{figure}[H]
      \centering
      \includegraphics[width=10cm]{screen-openproject-taskassigned.png}
      \caption{Tampilan Daftar \textit{Task} yang Diberikan kepada Penulis}
\end{figure}

\subsection{Objective Alignment}
Sebagai bagian dari penyesuaian terhadap sistem kerja perusahaan, mahasiswa
juga mengikuti kegiatan Objective Alignment yang dilaksanakan pada 4 Maret 2026
pukul 16.00–17.30 WIB di Ruang Auditorium. Kegiatan ini bertujuan untuk
menyelaraskan pemahaman seluruh anggota tim terhadap target kerja, prioritas
pengembangan, serta arah pencapaian yang ingin diwujudkan oleh divisi terkait.
Keikutsertaan mahasiswa dalam kegiatan ini memberikan pemahaman bahwa
pengembangan aplikasi tidak hanya berfokus pada aspek teknis, tetapi juga harus
sejalan dengan kebutuhan bisnis dan tujuan organisasi.

\begin{figure}[H]
      \centering
      \includegraphics[width=10cm]{event-objectivealignment.png}
      \caption{Kegiatan \textit{Objective Alignment}}
\end{figure}

